(1) 不可压Bingham流体位于两静止平板 \(y=\pm h\) 之间，恒定压强梯度 \(-dp/dx=G>0\)，塑性黏度 \(\eta>0\)、屈服应力 \(\tau_y>0\)。忽略入口效应，推导剪应力场，判定 \(Gh\leq\tau_y\) 时的唯一静止解，并确定 \(Gh>\tau_y\) 的未屈服核心半宽 \(a\)。
(2) 取 \(Gh>\tau_y\)，本构为 \(\tau_{xy}=\tau_y\operatorname{sgn}(u')+\eta u'\)（屈服区）。求 \(u(y)\) 的完整分段解，验证壁面无滑移、界面速度与剪切率连续。
(3) 求单位跨向宽度体积流率 \(Q=\int_{-h}^h u\,dy\) 的闭式，化成无量纲屈服比 \(a/h\) 的多项式，证明随 \(\tau_y\) 增大流率严格下降，并给出牛顿极限。
(4) 求单位长度、单位跨向宽度耗散 \(\mathcal D=\int_{-h}^h\tau_{xy}u'\,dy\)，从功率平衡独立证明 \(\mathcal D=GQ\)。解释塞状区为何不产生黏塑性耗散及临界屈服附近的流率标度。
